\documentclass[UTF8,11pt]{ctexart}
\usepackage[a4paper,margin=24mm]{geometry}
\usepackage{amsmath,amssymb,amsthm,mathtools,mathrsfs}
\usepackage[all]{xy}
\usepackage{xurl,hyperref,enumitem}
\usepackage{tikz}
\usetikzlibrary{arrows.meta,cd,decorations.markings,calc,patterns}
\hypersetup{hidelinks,breaklinks=true}
\setlength{\emergencystretch}{3em}
\allowdisplaybreaks
\newtheorem*{theorem}{Theorem}
\newtheorem*{thm}{Theorem}
\newtheorem*{lemma}{Lemma}
\newtheorem*{proposition}{Proposition}
\newtheorem*{prop}{Proposition}
\newtheorem*{corollary}{Corollary}
\newtheorem*{cor}{Corollary}
\newtheorem*{claim}{Claim}
\theoremstyle{definition}
\newtheorem*{definition}{Definition}
\newtheorem*{defn}{Definition}
\newtheorem*{remark}{Remark}
\newtheorem*{example}{Example}
\newcommand{\R}{\mathbb R}
\newcommand{\C}{\mathbb C}
\newcommand{\Z}{\mathbb Z}
\newcommand{\Q}{\mathbb Q}
\newcommand{\N}{\mathbb N}
\newcommand{\F}{\mathbb F}
\newcommand{\D}{\mathbb D}
\newcommand{\bB}{\mathbb B}
\newcommand{\bC}{\mathbb C}
\newcommand{\bR}{\mathbb R}
\newcommand{\bZ}{\mathbb Z}
\newcommand{\bN}{\mathbb N}
\newcommand{\bQ}{\mathbb Q}
\newcommand{\bD}{\mathbb D}
\newcommand{\bF}{\mathbb F}
\newcommand{\CP}{\mathbb{CP}}
\newcommand{\RP}{\mathbb{RP}}
\newcommand{\sO}{\mathscr O}
\newcommand{\sI}{\mathscr I}
\newcommand{\sF}{\mathscr F}
\newcommand{\sG}{\mathscr G}
\newcommand{\sH}{\mathscr H}
\newcommand{\sR}{\mathscr R}
\newcommand{\sM}{\mathscr M}
\newcommand{\sV}{\mathscr V}
\newcommand{\sU}{\mathscr U}
\DeclareMathOperator{\Hom}{Hom}
\DeclareMathOperator{\Ext}{Ext}
\DeclareMathOperator{\Tor}{Tor}
\DeclareMathOperator{\Spec}{Spec}
\DeclareMathOperator{\Proj}{Proj}
\DeclareMathOperator{\supp}{supp}
\DeclareMathOperator{\im}{im}
\DeclareMathOperator{\id}{id}
\DeclareMathOperator{\Tr}{Tr}
\DeclareMathOperator{\tr}{tr}
\DeclareMathOperator{\rank}{rank}
\DeclareMathOperator{\ord}{ord}
\DeclareMathOperator{\dist}{dist}
\DeclareMathOperator{\End}{End}
\DeclareMathOperator{\Aut}{Aut}
\DeclareMathOperator{\Gal}{Gal}
\DeclareMathOperator{\vol}{vol}
\DeclareMathOperator{\Cl}{Cl}
\DeclareMathOperator{\coker}{coker}
\DeclareMathOperator{\codim}{codim}
\DeclareMathOperator{\divergence}{div}
\DeclareMathOperator{\Ric}{Ric}
\DeclareMathOperator{\grad}{grad}
\DeclareMathOperator{\Hess}{Hess}
\newcommand{\abs}[1]{\left\lvert#1\right\rvert}
\newcommand{\sabs}[1]{\lvert#1\rvert}
\newcommand{\babs}[1]{\bigl\lvert#1\bigr\rvert}
\newcommand{\Babs}[1]{\Bigl\lvert#1\Bigr\rvert}
\newcommand{\bbabs}[1]{\biggl\lvert#1\biggr\rvert}
\newcommand{\BBabs}[1]{\Biggl\lvert#1\Biggr\rvert}
\newcommand{\norm}[1]{\left\lVert#1\right\rVert}
\newcommand{\snorm}[1]{\lVert#1\rVert}
\newcommand{\bnorm}[1]{\bigl\lVert#1\bigr\rVert}
\newcommand{\Bnorm}[1]{\Bigl\lVert#1\Bigr\rVert}
\newcommand{\bbnorm}[1]{\biggl\lVert#1\biggr\rVert}
\newcommand{\BBnorm}[1]{\Biggl\lVert#1\Biggr\rVert}
\newcommand{\widebar}[1]{\overline{#1}}
\newcommand{\nicefrac}[2]{\frac{#1}{#2}}
\newcommand{\linnprod}[2]{\langle#1,#2\rangle}
\newcommand{\myindex}[1]{#1}
\newcommand{\glsadd}[1]{}
\newcommand{\avoidbreak}{}
\newcommand{\sourceref}[1]{\textnormal{[原文：\nolinkurl{#1}]}}
\newcommand{\thmref}[1]{\sourceref{#1}}
\newcommand{\lemmaref}[1]{\sourceref{#1}}
\newcommand{\propref}[1]{\sourceref{#1}}
\newcommand{\corref}[1]{\sourceref{#1}}
\newcommand{\defnref}[1]{\sourceref{#1}}
\newcommand{\exampleref}[1]{\sourceref{#1}}
\newcommand{\exerciseref}[1]{\sourceref{#1}}
\newcommand{\figureref}[1]{\sourceref{#1}}
\newcommand{\sectionref}[1]{\sourceref{#1}}
\newcommand{\appendixref}[1]{\sourceref{#1}}
\newcommand{\stacksref}[2]{\href{https://stacks.math.columbia.edu/tag/#1}{#2 [#1]}}


\newcommand{\E}{\mathbb E}
\newcommand{\Prob}{\mathbb P}
\newcommand{\Reff}{\mathcal R_{\mathrm{eff}}}
\newcommand{\eps}{\varepsilon}
\DeclareMathOperator{\diam}{diam}
\DeclareMathOperator{\Var}{Var}
\DeclareMathOperator{\Crit}{Crit}
\newcommand{\dd}{\,\mathrm d}


\begin{document}
\section*{084\quad 任意空间对的CW逼近构造与相对同伦唯一性}
\subsection*{来源与整理范围}
Bernard Badzioch，MTH 727 Algebraic Topology，Chapter 15, Weak Homotopy Type；原标签 CW APPROX DEF、CW APPROX THM 及其全部证明。
\par\url{https://github.com/bbadzioch/topology_lecture_notes/blob/master/MTH_727_notes/15_weak_homotopy_type.tex}
\par\textbf{恢复方式（整理者说明）：}依据人类原文重新整理以补齐缺失题号；并非旧题证文件的逐字节恢复；2026-09-28。
\subsection*{作者命题与人类证明}

\begin{definition}[CW approximation]
A CW approximation of a space $X$ is a CW complex $Y$ together with a weak
homotopy equivalence $f:Y\to X$. More generally, a CW approximation of a pair
$(X,A)$ is a relative CW complex $(Y,A)$ and a weak equivalence $f:Y\to X$
such that $f|_A=\id_A$.
\end{definition}
\begin{theorem}[CW approximation; original label CW APPROX THM]
Any pair $(X,A)$ has a CW approximation. Any two CW approximations of this pair
are homotopy equivalent relative to $A$.
\end{theorem}
\begin{proof}
Assume first that $X$ is path connected. We construct relative CW complexes
$(Y^{(n)},A)$ and compatible maps $f^{(n)}:Y^{(n)}\to X$ so that
$Y^{(n)}$ is obtained from $Y^{(n-1)}$ by attaching $n$-cells, the maps restrict
to the identity on $A$, and
\[
 f^{(n)}_*:\pi_i(Y^{(n)})\longrightarrow\pi_i(X)
\]
is an isomorphism for $i<n$ and an epimorphism for $i=n$.
The union of the maps will give the required CW approximation.

Let $\{A_i\}_{i\in I}$ be the path components of $A$, choose $a_i\in A_i$,
and fix $x_0\in X$. Add a new $0$-cell $e^0$ to $A$. For each $i$, attach a
$1$-cell joining $e^0$ to $a_i$. For each class
$[\tau:(S^1,s_0)\to(X,x_0)]\in\pi_1(X,x_0)$, attach a circle $S^1_\tau$
by identifying $s_0$ with $e^0$. Call the resulting space $Y^{(1)}$.
Since $X$ is path connected, choose a path $\omega_i$ from $x_0$ to $a_i$.
Define $f^{(1)}$ to be the identity on $A$, to send $e^0$ to $x_0$, to map
the joining edge using $\omega_i$, and to map $S^1_\tau$ using $\tau$.
It induces a bijection on $\pi_0$ and a surjection on $\pi_1$.

Suppose now that the construction has been performed through dimension $n$.
For every element
\[
 [\omega:(S^n,s_0)\to(Y^{(n)},e^0)]\in\ker f^{(n)}_*,
\]
attach an $(n+1)$-cell using $\omega$ as attaching map. Denote the result by
$\overline Y^{(n+1)}$. Since $[f^{(n)}\omega]=0$, the composite on the
boundary extends over $D^{n+1}$, giving an extension
$\overline f^{(n+1)}:\overline Y^{(n+1)}\to X$.
Next, for each $[\tau]\in\pi_{n+1}(X,x_0)$ attach a sphere
$S^{n+1}_\tau$ at $e^0$, and extend the map by $\tau$ on that sphere.
This gives $Y^{(n+1)}$ and $f^{(n+1)}$.

With $i:Y^{(n)}\hookrightarrow Y^{(n+1)}$ the inclusion, we have
\[
\begin{tikzcd}[column sep=large]
\pi_n(Y^{(n)},e^0) \arrow[r,"i_*"] \arrow[dr,"f^{(n)}_*"'] &
\pi_n(Y^{(n+1)},e^0) \arrow[d,"f^{(n+1)}_*"]\\
&\pi_n(X,x_0).
\end{tikzcd}
\]
Surjectivity of $f^{(n)}_*$ gives surjectivity of $f^{(n+1)}_*$ in dimension
$n$, and the cells attached along its kernel make this map injective.
The lower homotopy groups are unchanged. The newly attached spheres make
$f^{(n+1)}_*$ surjective in dimension $n+1$. Thus the inductive requirements
hold. Taking $Y=\bigcup_nY^{(n)}$ and $f=\bigcup_nf^{(n)}$ gives a weak
equivalence $Y\to X$ and a relative CW structure on $(Y,A)$.

If $X$ is not path connected, let $X_i$ be its path components and apply the
construction to each pair $(X_i,A\cap X_i)$. A CW approximation of $(X,A)$ is
obtained from
\[
 A\sqcup\coprod_iY_i
\]
by identifying each point of $A\cap X_i\subset Y_i$ with the corresponding
point of $A$.

For uniqueness, suppose $f_i:(Y_i,A)\to(X,A)$ are CW approximations,
$i=1,2$. The relative homotopy lifting property for a weak equivalence and
a relative CW complex (original Corollary HELP COR) gives a map
$g:Y_1\to Y_2$ restricting to the identity on $A$ and satisfying
$f_2g\simeq f_1$ relative to $A$. Similarly, there is $h:Y_2\to Y_1$
with $f_1h\simeq f_2$ relative to $A$. Consequently there is a homotopy
$\varphi:Y_1\times I\to X$ from $f_1$ to $f_1hg$, relative to $A$.

Set
\[
 Z_1=Y_1\times\{0,1\}\cup A\times I\subset Y_1\times I.
\]
The pair $(Y_1\times I,Z_1)$ is a relative CW complex. Define
$\psi:Z_1\to Y_1$ by
\[
 \psi(y,t)=\begin{cases}y,&t<1,\\hg(y),&t=1.\end{cases}
\]
These prescriptions agree on $A\times\{1\}$ because $hg|_A=\id_A$.
We have the commutative diagram
\[
\begin{tikzcd}[column sep=large]
Z_1 \arrow[r,"\psi"] \arrow[d,hook] &Y_1\arrow[d,"f_1"]\\
Y_1\times I \arrow[r,"\varphi"'] \arrow[ur,dashed,"\overline\varphi"] & X.
\end{tikzcd}
\]
Applying the same lifting property gives $\overline\varphi:Y_1\times I\to Y_1$
extending $\psi$. It is a homotopy $\id_{Y_1}\simeq hg$ relative to $A$.
The analogous construction gives $\id_{Y_2}\simeq gh$ relative to $A$.
\end{proof}

\subsection*{整理说明与先修范围（非作者正文）}
依据已取得的原生 TeX 逐段转录，统一原文的缩写宏与图式排版；不是上游文件的逐字节副本。原文末段将提升映射的值域写成 $X$，此处按该段图式与同伦端点改为 $Y_1$；核的记号补上诱导同态星号。保留不连通空间和相对唯一性，未将其删成只证明存在。

允许引用相对同伦群、胞腔附着对低维同伦群的影响、CW 弱拓扑与紧致域映射性质，以及原 Chapter 14 的相对同伦提升推论：弱等价 $E\to B$ 对相对 CW 对的可交换图存在相对同伦提升。该推论在此作为明确的标准先修。难点是按核与像交替附着胞腔，并将比较映射的相容性转化为保持 $A$ 的全局同伦；并非直接引用 CW 逼近定理。
\subsection*{可修改的简短标注（非作者正文）}
$c$：一般拓扑空间的胞腔模型与相对同伦类型

$a$：同伦群的核；胞腔附着；球面代表；相对CW结构；弱等价的同伦提升

\texttt{cross\_theory\_status = yes}。把同伦群中缺少的代表和多余的核分别转化为球面、胞腔附着，再用提升图式把代数上的弱等价转回相对同伦等价；见归纳构造及最后两个图式。

全部标注为非穷尽、可修改初稿。
\subsection*{来源许可与编辑归属}
作者公开讲义按 CC BY-NC-SA 4.0 使用；保留作者与原仓库归属，编辑转录沿用相同许可。
ChatGPT 加入题号、中文来源说明、整理范围和初稿标注；编辑说明与人类证明分列。
\end{document}
